\section{One Creator, many servants}
\label{sec:faith}

An angel's invisibility does not place it halfway between God and creation.
The decisive distinction is between the One who gives being and everything
that receives it. The Fourth Lateran Council's profession of faith names
one Creator of spiritual and bodily creatures, explicitly including the
angelic creation. It also declares that the devil and the other demons
were created good and became evil through their own action.\footnote{Lateran IV
(1215), constitution~1, \emph{Firmiter credimus}, DS~800--801;
Latin text in the identified Denzinger web witness.}

\subsection{Creation is a gift, not a division of God}

Vatican I repeats Lateran's profession in \emph{Dei Filius} and explains
that God creates freely, communicating goods without increasing his own
happiness. Its canons reject the claim that spiritual things emerge from
the divine substance. An angel is therefore neither a detached portion of
God nor a necessary overflow that God could not withhold. Creation means
that the creature's entire being is received.\footnote{Vatican I,
\emph{Dei Filius} (24 April 1870), ch.~1; canons on God the Creator,
1--5, particularly 4--5. Latin text, Holy See.}

This confession explains why the greatest angel and the smallest bodily
creature share something fundamental: neither exists of itself. Their
perfections differ immensely, but both depend on the Creator. Aquinas's
account of separate substances supplies a philosophical articulation:
freedom from matter does not make an angel identical with its own act of
existing. Immateriality answers one question about composition; dependence
on God answers another.\footnote{\ST{I q.~50 a.~2 ad~3}; compare
\emph{De ente et essentia}, chapter beginning \emph{Nunc restat videre}, in the Baur--Koch Latin witness.}

The council's confession also locates evil in a misuse of created goodness.
The fallen spirit is not an uncreated rival to God. Its intelligence and
power remain finite creaturely goods, while the disorder of its will is
its own. This makes moral evil intelligible without attributing an evil
nature to the Creator.\footnote{Lateran IV, \emph{Firmiter}, DS~800;
\ST{I q.~63 aa.~1--4}.}

\subsection{Personal spirits and the freedom of grace}

The \emph{Catechism of the Catholic Church} teaches that angels are
spiritual, personal, intelligent and willing creatures; their service is
centered on Christ and extends to the Church and human life. Its teaching
also holds together the fallen angels' original goodness, their free and
irrevocable rejection of God, and the limits of their power under
providence.\footnote{\emph{Catechism of the Catholic Church},
328--336, 391--395, official English web witness consulted 29 September
2026. These summaries are \textbf{Church teaching}; the original goodness
and voluntary fall of the demons also have the conciliar basis above.}

Personhood matters here. An angel can know and love; it is not simply a
name for an impersonal force. Pius XII's \emph{Humani generis} rejects
doubts about angelic personality and the essential distinction between
matter and spirit. The same paragraph defends the gratuity of the
supernatural order against the claim that God must call every intellectual
creature to the beatific vision.\footnote{Pius XII, \emph{Humani generis}
(12 August 1950), no.~26.}

The natural greatness of an intellect consequently does not earn it the
vision of God's essence. In Aquinas's account, even a creature that knows
without bodily senses must be elevated by grace. The difference between
nature and glory will govern the treatment of angelic knowledge, love,
merit, and perseverance. The most exalted created understanding still
receives its beatitude.\footnote{\ST{I q.~56 a.~3; q.~62 aa.~1--2}.}

\subsection{The abundance of created goodness}

The multitude of angels expresses the abundance of divine generosity.
Aquinas locates their perfection within the good of the whole creation:
God's goodness is represented through many beings and many kinds of
perfection. The angels have a place in this order before they perform
any particular service for us. Their ministry flows from a goodness
already received.\footnote{\ST{I q.~50 aa.~1, 3}.}

The same generosity makes the good of one creature fruitful for another.
The intellect that receives illumination communicates it; the spirit
established in blessedness serves a human pilgrim; the heavenly host
shares its praise with the Church. Created excellence reaches its proper
expression in communion. The Fathers unfold this order as a work of
divine wisdom, holiness and love.\footnote{\ST{I q.~106 aa.~1, 4;
q.~112 a.~1; q.~113 a.~1}; Hebrews 12:22--24.}
